% !PRINTED 328
\weberpage{328}
%<<CONT
is satisfied. If in particular the form is independent of the variables $x$ and depends only on the coefficients $a$, it becomes an \gesperrt{invariant}. The exponent $\lambda$, which is called the \gesperrt{weight} of the covariant, is an integer.

In Vol.\ I, \S.\,59 we became acquainted with a covariant of weight 2 that is present in all forms whose degree is greater than 1, namely the \gesperrt{Hesse determinant}:

\weq{2}{H(x_1, x_2, x_3) = \begin{vmatrix} f''(x_1, x_1), & f''(x_1, x_2), & f''(x_1, x_3) \\ f''(x_2, x_1), & f''(x_2, x_2), & f''(x_2, x_3) \\ f''(x_3, x_1), & f''(x_3, x_2), & f''(x_3, x_3) \end{vmatrix}.}

Besides this covariant $H$ we shall consider two other general covariants, which exist for all ternary forms of degree higher than $2$, and which can also easily be generalised to forms of arbitrarily many variables. The first of these forms, if $H$ denotes the \gesperrt{Hesse} covariant as above, is:

\weq{3}{C_1(x, a) = \begin{vmatrix} f''(x_1, x_1), & f''(x_1, x_2), & f''(x_1, x_3), & H'(x_1) \\ f''(x_2, x_1), & f''(x_2, x_2), & f''(x_2, x_3), & H'(x_2) \\ f''(x_3, x_1), & f''(x_3, x_2), & f''(x_3, x_3), & H'(x_3) \\ H'(x_1), & H'(x_2), & H'(x_3) & 0 \end{vmatrix}.}

To demonstrate the invariant property for this form, we make in it the substitution \S.\,85, (1). Thereby $f(x, a)$ passes over into $\varphi(y, b)$ and $r^2 H(x, a)$ into $H(y, b)$. If we then understand by $H'(y_1), H'(y_2), H'(y_3)$ the derivatives of $H(y, b)$, we have the formulae

\weq{4}{
\begin{aligned}
\varphi''(y_\mu, y_\nu) &= \sum\limits^{i,k}_{i} f''(x_i, x_k)\, \alpha_i^{(\mu)} \alpha_k^{(\nu)} \\
H'(y_\nu) &= r^2 \sum H'(x_i)\, \alpha_i^{(\nu)},
\end{aligned}
}
in which $\mu, \nu, i, k$ run from 1 to 3.

If one now multiplies the determinant $C_1(x, a)$ twice in succession by the substitution determinant, written in the form

\[
r = \begin{vmatrix} \alpha_1^{(1)}, & \alpha_1^{(2)}, & \alpha_1^{(3)}, & 0 \\ \alpha_2^{(1)}, & \alpha_2^{(2)}, & \alpha_2^{(3)}, & 0 \\ \alpha_3^{(1)}, & \alpha_3^{(2)}, & \alpha_3^{(3)}, & 0 \\ 0, & 0, & 0, & 1 \end{vmatrix}
\]
by taking the one
%CONT>>
